\exercisegroup

\begin{enumerate}
\def\labelenumi{\arabic{enumi})}
\item
  设 G 是拓扑群 O(2, R)。有 \(\pi_{0}(G) \simeq Z/2Z\) 和 \(\pi_{1}(G) \simeq Z\)；\(\pi_{0}(G)\) 的非平凡元作用于 \(\pi_{1}(G)\)，其映射为 \(t \mapsto -t\)。
\item
  记 H 为哈密顿四元数代数（TG，第 VIII 章，第 4 页），\((1, i, j, k)\) 为其典范基。记 N 为 H 的范数二次型；范数为 1 的四元数集是球面 \(S_{3}\)。令 \(H_{0}\) 为迹为零的四元数所成的实向量子空间，令 \(N_{0}\) 为 N 在 \(H_{0}\) 上的限制。
\end{enumerate}

\begin{enumerate}
\def\labelenumi{\alph{enumi})}
\item
  证明 TG，第 VIII 章，第 4 页，命题 4 中的同态使球面 \(\mathbf{S}_3\) 成为 \(\mathbf{SO}(\mathrm{N}_0)\) 的 2 次覆叠。
\item
  由此推出 SO(3, R) 的庞加莱群同构于 Z/2Z。具体描述 SO(3, R) 中 \(I_{3}\) 内的一个圈，其严格同伦类是该庞加莱群的唯一非平凡元。
\item
  设 \(\varphi\) 是从 \(\mathbf{S}_3 \times \mathbf{S}_3\) 到 SO(4, R) 的映射，由 \(\varphi(\mathbf{q}_1, \mathbf{q}_2)(\mathbf{x}) = \mathbf{q}_1 \mathbf{x} \mathbf{q}_2^{-1}\) 给出。证明 \(\varphi\) 使 \(\mathbf{S}_3 \times \mathbf{S}_3\) 成为 2 次覆叠。
\item
  推出 SO(4, R) 的庞加莱群同构于 Z/2Z。
\end{enumerate}

\begin{enumerate}
\def\labelenumi{\arabic{enumi})}
\setcounter{enumi}{2}
\tightlist
\item
  设 n 是自然数。令拓扑群 \(\mathrm{SO}(n+1,\mathbf{R})\) 作用于 \(S_{n}\)，即 \(R^{n+1}\) 的球面。
\end{enumerate}

\begin{enumerate}
\def\labelenumi{\alph{enumi})}
\item
  证明向量 \(e_{n+1}\) 的固定子可识别为拓扑群 \(\mathrm{SO}(n,\mathbf{R})\)。
\item
  设 \(f \colon \mathbf{I} \to \mathbf{S}_n\) 是连续映射。假设 \(f(t) \neq (0, \ldots, 0, 1)\) 对所有 \(t \in \mathbf{I} .\) 都成立。证明，存在连续映射 \(g \colon \mathbf{I} \to \mathrm{SO}(n + 1, \mathbf{R})\)，使得 \(f(t) = g(t) \cdot \mathbf{e}_{n+1}\) 对所有 \(t \in \mathbf{I} \)。
\item
  若 \(n \geqslant 3\)，证明 \(\mathrm{SO}(n, \mathbf{R})\) 到 \(\mathrm{SO}(n+1, \mathbf{R})\) 的包含映射在庞加莱群之间诱导同构。
\end{enumerate}

\begin{enumerate}
\def\labelenumi{\arabic{enumi})}
\setcounter{enumi}{3}
\tightlist
\item
  设 n 是自然数。令 \(B_{0}\) 为 \(\mathbf{SL}(n,\mathbf{R})\) 中由对角系数全为严格正数的上三角矩阵组成的子群。
\end{enumerate}

\begin{enumerate}
\def\labelenumi{\alph{enumi})}
\item
  证明从 \(\mathrm{SO}(n,\mathbf{R})\times\mathrm{B}_{0}\) 到 \(\mathbf{SL}(n,\mathbf{R})\) 的映射 \((g,u)\mapsto g\cdot u\) 是同胚（参见 INT，第 VII 章，第 91 页，命题 7）。
\item
  证明 SO(n, R) 到 SL(n, R) 的包含映射在庞加莱群之间诱导同构。
\end{enumerate}

\begin{enumerate}
\def\labelenumi{\arabic{enumi})}
\setcounter{enumi}{4}
\tightlist
\item
  设 G 是连通的可解圈拓扑群，设 \((\widetilde{\mathrm{G}}, p)\) 是 G 的万有覆叠；在 \(\widetilde{G}\) 上赋予拓扑群结构，使 p 成为群同态。记 Z 为 G 的中心，e 为其单位元，\(\widetilde{Z}\) 为 \(\widetilde{G}\) 的中心，\(\widetilde{e}\) 为其单位元。
\end{enumerate}

\begin{enumerate}
\def\labelenumi{\alph{enumi})}
\item
  设 \(f\) 是群 G 的连续自同构。证明，存在唯一的连续映射 \(\widetilde{f} \colon \widetilde{\mathrm{G}} \to \widetilde{\mathrm{G}}\)，使得 \(\widetilde{f}(\widetilde{e}) = \widetilde{e}\) 且 \(p \circ \widetilde{f} = f \circ p\)。证明 \(\widetilde{f}\) 是群自同构。证明映射 \(\varphi \colon f \mapsto \widetilde{f}\) 是从 \(\operatorname{Aut}(\mathbf{G})\) 到 \(\operatorname{Aut}(\widetilde{\mathbf{G}})\) 的单射群同态。
\item
  记 \(\operatorname{Out}(\mathbf{G})\) 为群 \(\operatorname{Aut}(\mathbf{G})\) 对内自同构子群的商；类似地定义 \(\operatorname{Out}(\widetilde{\mathbf{G}})\)。证明同态 \(\varphi\) 通过取商所诱导的同态 \(\psi\colon\operatorname{Out}(\mathbf{G})\to\operatorname{Out}(\widetilde{\mathbf{G}})\) 是单射群同态。
\end{enumerate}

\begin{enumerate}
\def\labelenumi{\arabic{enumi})}
\setcounter{enumi}{5}
\tightlist
\item
  设 G、H 是连通实李群，且 \(f: G \to H\) 是李群的满射态射。记 \(Z_{G}\) 为 G 的中心，\(Z_{H}\) 为 H 的中心。
\end{enumerate}

\begin{enumerate}
\def\labelenumi{\alph{enumi})}
\item
  证明以下断言等价：(i) 映射 f 使 G 成为 H 的覆叠；(ii) \(\operatorname{Ker}(f)\) 是离散的；(iii) f 通过取李代数诱导的态射 \(\mathrm{L}(f)\colon\mathrm{L}(\mathrm{G})\to\mathrm{L}(\mathrm{H})\) 是同构。
\item
  证明此时 \(f(\mathrm{Z}_{\mathrm{G}})=\mathrm{Z}_{\mathrm{H}}\)，且 \(f^{-1}(\mathrm{Z}_{\mathrm{H}})=\mathrm{Z}_{\mathrm{G}}\)。
\end{enumerate}

\begin{enumerate}
\def\labelenumi{\arabic{enumi})}
\setcounter{enumi}{6}
\tightlist
\item
  设 G 是实李群，记 \(G_{0}\) 为其单位分支，\(Z_{0}\) 为 \(G_{0}\) 的中心，\(\pi_{0}(G)=G/G_{0}\)。设 \((\widetilde{G}_{0},p_{0})\) 是 \(G_{0}\) 的万有覆叠，\(\widetilde{Z}_{0}\) 是 \(\widetilde{G}_{0}\) 的中心，\(\widetilde{e}\) 是其单位元。
\end{enumerate}

\begin{enumerate}
\def\labelenumi{\alph{enumi})}
\item
  设 \(\theta\colon\pi_{0}(\mathrm{G})\to\mathrm{Out}(\mathrm{G}_{0})\) 是与扩张 \(\mathcal{E}\colon\mathrm{G}_{0}\to\mathrm{G}\to\pi_{0}(\mathrm{G})\) 相关联的同态。证明文献所引处定义的上同调类 \(\omega(\pi_{0}(\mathrm{G}),\mathrm{Z}_{0},\theta)\)（它属于 \(\mathrm{H}^{3}(\pi_{0}(\mathrm{G}),\mathrm{Z}_{0})\)）为零。
\item
  假设存在李群 \(\widetilde{\mathrm{G}}\)、满射 \(p\colon\widetilde{\mathrm{G}}\to\mathrm{G}\)，以及同构 \(j\)，它把 \(\widetilde{\mathrm{G}}_{0}\) 同构到 \(\widetilde{\mathrm{G}}\) 的单位分支上，使得 \(p_{0}=p\circ j\)。证明 \(\omega(\pi_{0}(\mathrm{G}),\widetilde{\mathrm{Z}}_{0},\psi\circ\theta)=0\) 在 \(\mathrm{H}^{3}(\pi_{0}(\mathrm{G}),\widetilde{\mathrm{Z}}_{0})\) 中，其中 \(\psi\colon\mathrm{Out}(\mathrm{G}_{0})\to\mathrm{Out}(\widetilde{\mathrm{G}}_{0})\) 是练习 5 中定义的同态。
\item
  假设 \(\omega(\pi_{0}(\mathrm{G}),\widetilde{\mathrm{Z}}_{0},\psi\circ\theta)=0\)。设 \(\widetilde{E}\colon\widetilde{G}_{0}\to E\to\pi_{0}(G)\) 是以 \(\pi_{0}(G)\) 为商、以 \(\widetilde{G}_{0}\) 为核的扩张，且其相关联同态等于 \(\psi\circ\theta\)。证明 \(\mathrm{E}/\pi_{1}(\mathrm{G}_{0})\) 是以 \(\pi_{0}(G)\) 为商、以 \(G_{0}\) 为核的扩张，且其相关联同态等于 \(\theta\)。
\end{enumerate}

现在令 \(\gamma\) 为 \(\mathrm{H}^2 (\pi_0(\mathrm{G}),\mathrm{G}_0)\) 中满足 \(\gamma \cdot [\mathrm{E} / \pi_1(\mathrm{G}_0)] = [\mathrm{G}]\) 的唯一元素（所引文献，e)）。考虑同态 \(\delta \colon \mathrm{H}^2 (\pi_0(\mathrm{G}),\mathrm{Z}_0)\to\) \(\mathrm{H}^3 (\pi_0(\mathrm{G}),\pi_1(\mathrm{G}))\)，它由阿贝尔群正合序列 \(1\to \pi_1(\mathrm{G}_0)\to\) \(\widetilde{\mathbf{Z}}_0\to \mathbf{Z}_0\to 1\)（练习 6）关联。证明元素 \(\delta (\gamma)\)（属于 \(\mathrm{H}^3 (\pi_0(\mathrm{G}),\pi_1(\mathrm{G}))\)）不依赖于扩张 \(\widetilde{\mathcal{E}}\) 的选择；将其记为 \(\delta (\mathrm{G})\)。

d) 假设 \(\omega (\pi_0(\mathrm{G}),\widetilde{\mathrm{Z}}_0,\psi \circ \theta) = 0\)。存在李群 \(\widetilde{\mathrm{G}}\)、满射 \(p\colon \widetilde{\mathrm{G}}\to \mathrm{G}\) 以及同构 \(j\)，它把 \(\widetilde{\mathrm{G}}_0\) 同构到 \(\widetilde{G}\) 的单位分支上，使得 \(p_{0}=p\circ j\) 的必要且充分条件是 \(\delta(\mathrm{G})=0.\)（5）

(5) 例子参见 R. L. TAYLOR 的文章《非连通拓扑群的覆叠群》，《美国数学学会会刊》5 (1954)，第 753–768 页。

\begin{enumerate}
\def\labelenumi{\arabic{enumi})}
\setcounter{enumi}{7}
\tightlist
\item
  设 G 是李群，H 是 G 的闭子群；记 \(H_{0}\) 为 H 的单位分支，\(\pi_{0}(H)\) 为群 \(H/H_{0}\)。再记 p 为从 G 到 G/H 的典范满射。
\end{enumerate}

\begin{enumerate}
\def\labelenumi{\alph{enumi})}
\item
  映射 p 具有道路提升性质。
\item
  若 H 连通，则同态 \(p_{*}\colon\pi_{1}(\mathrm{G},e)\to\pi_{1}(\mathrm{G}/\mathrm{H},p(e))\) 是满射。
\item
  存在唯一映射 \(\varphi\colon\pi_{1}(\mathrm{G/H},p(e))\to\pi_{0}(\mathrm{H})\)，对于 G 中每条道路 \(\widetilde{c}\)，其起点为 \(e\) 且终点属于 H，它把 G/H 中圈 \(p\circ\widetilde{c}\) 的道路类对应到 \(\widetilde{c}(1)\) 模 \(\mathrm{H}_{0}\) 的类；它是群同态。
\item
  若 G 单连通，则 \(\varphi\) 是同构。特别地，H 连通当且仅当 G/H 单连通。
\end{enumerate}

\begin{enumerate}
\def\labelenumi{\arabic{enumi})}
\setcounter{enumi}{8}
\tightlist
\item
  设 G 是连通李群，\((\widetilde{\mathbf{G}}, q)\) 是 G 的万有覆叠，\(\widetilde{e}\) 是 \(\widetilde{G}\) 的单位元。设 H 是 G 的闭连通子群；记 \(i: H \to G\) 为典范注入，\(p: G \to G/H\) 为典范满射。
\end{enumerate}

\begin{enumerate}
\def\labelenumi{\alph{enumi})}
\item
  令 \(H_{1}\) 为 \(\overline{q}^{-1}(H)\) 的单位分支，令 \(q_{1}\colon\widetilde{\mathrm{G}}/\mathrm{H}_{1}\to\mathrm{G}/\mathrm{H}\) 为 q 通过取商诱导的映射。证明 \(q_{1}\) 使 \(\widetilde{G}/H_{1}\) 成为 G/H 的覆叠，且其在 \(p(e)\) 上的纤维同构于 \(\pi_{1}(\mathrm{G}/\mathrm{H},p(e))\)。
\item
  推出存在群的正合序列
\end{enumerate}

\[
\pi_ {1} (\mathrm{H} _ {1}, \widetilde {e}) \xrightarrow {q _ {*}} \pi_ {1} (\mathrm{H}, e) \xrightarrow {i _ {*}} \pi_ {1} (\mathrm{G}, e) \xrightarrow {p _ {*}} \pi_ {1} (\mathrm{G} / \mathrm{H}, p (e)).
\]

\begin{enumerate}
\def\labelenumi{\arabic{enumi})}
\setcounter{enumi}{9}
\tightlist
\item
  \begin{enumerate}
  \def\labelenumii{\alph{enumii})}
  \tightlist
  \item
    证明，对于每个整数 \(n \geqslant 2\)，群 \(\mathrm{SU}(n, \mathbf{C})\) 都是单连通的。（空间 \(\mathrm{SU}(2, \mathbf{C})\) 同胚于 \(S_{3}\)。令 \(\mathrm{SU}(n, \mathbf{C})\) 作用于 \(C^{n}\) 的单位球面，把点 \((0, \ldots, 0, 1)\) 的稳定子识别为群 \(\mathrm{SU}(n - 1, \mathbf{C})\)，并使用练习 9。）
  \end{enumerate}
\end{enumerate}

\begin{enumerate}
\def\labelenumi{\alph{enumi})}
\setcounter{enumi}{1}
\tightlist
\item
  证明，对于每个整数 \(n \geqslant 2\)，群 \(\mathbf{SL}(n, \mathbf{C})\) 都是单连通的。（令 \(B_{+}\) 为 \(\mathbf{SL}(n, \mathbf{C})\) 中由对角系数为实数 \textgreater{} 0 的上三角矩阵组成的子群。证明从 \(\mathrm{SU}(n, \mathbf{C}) \times \mathrm{B}_{+}\) 到 \(\mathbf{SL}(n, \mathbf{C})\) 的映射 \((g, u) \mapsto g \cdot u\) 是同胚。）
\end{enumerate}

